\documentclass[11pt]{article}
\usepackage[margin=1in]{geometry}
\usepackage{amsmath,amssymb,amsthm}
\usepackage[hidelinks]{hyperref}

\newtheorem{theorem}{Theorem}
\newtheorem{lemma}{Lemma}
\newtheorem{proposition}{Proposition}
\newtheorem{corollary}{Corollary}
\newtheorem{definition}{Definition}
\newtheorem{remark}{Remark}

\newcommand{\eps}{\varepsilon}
\newcommand{\N}{\mathbb{N}}
\newcommand{\Q}{\mathbb{Q}}
\newcommand{\R}{\mathbb{R}}
\newcommand{\Span}{\operatorname{span}}
\newcommand{\osc}{\operatorname{osc}}
\newcommand{\tp}{\operatorname{tp}}
\newcommand{\T}{\mathcal{T}}
\newcommand{\C}{\mathcal{C}}
\newcommand{\Sph}{\mathbb{S}}

\title{A Formula-Aware Krivine--Maurey Proof of\\Stability Implies Oscillation Stability}
\author{fa\_banach\_001, agent\_lane\_19}
\date{2026-06-16}

\begin{document}
\maketitle

\section*{Source Question}

James Hanson asks in \emph{Indiscernible Subspaces and Minimal Wide Types},
arXiv:2004.03062, Question 4.3:
\begin{quote}
Does (model theoretic) stability imply oscillation stability? That is to say,
if \(T\) is a stable Banach theory, is every unary formula oscillation stable
on models of \(T\)?
\end{quote}

This packet gives a claimed positive answer.

\section*{Proof Intuition}

The Dvoretzky--Milman route naturally produces finite-dimensional Hilbert
pieces on which a formula is nearly constant, but those pieces need not be
nested. That is why a direct compactness pullback from an ultrapower is too
weak: it would falsely prove oscillation stability for arbitrary uniformly
continuous Hilbert-sphere colorings.

The proof below uses a different stable-Banach mechanism. Krivine and Maurey,
in Raynaud's type-space formulation, prove that a stable Banach space contains
an \(\ell_p\)-type: in a suitable compact type semigroup, one has
\[
  c_1\tau * \cdots * c_n\tau = \|(c_1,\ldots,c_n)\|_p\,\tau .
\]
The key observation is that this argument can be run in an enriched type
space which records not only norm-type coordinates, but also all stable
formula coordinates relevant to a fixed unary instance \(\varphi(x,a)\). The
Raynaud/Krivine--Maurey proof uses only compactness, stability of coordinates,
dilation, convolution, minimal closed conic classes, and continuity points.
Those hypotheses remain true after adding stable formula coordinates.

Therefore the displayed equality is equality of enriched types. In particular,
if \(\|(c_i)\|_p=1\), then the coordinate recording the normalized pullback
\(\varphi(\sigma_c(\bar x),a)\) has the same value as the one-vector
coordinate. A block extraction converts the type into an actual basic sequence
inside the original subspace. The closed span of that block sequence is the
required oscillation-stable subspace.

\section*{Definitions}

Let \(M\models T\) be a Banach structure and let \(Y\subseteq M\) be an
infinite-dimensional closed linear subspace. A bounded uniformly continuous
function \(f\) on the unit sphere \(S_Y\) is oscillation stable on \(Y\) if for
every \(\eps>0\) there is an infinite-dimensional closed subspace
\(Z\subseteq Y\) such that
\[
  \sup\{|f(u)-f(v)|:u,v\in S_Z\}<\eps .
\]
A unary formula \(\varphi(x,a)\) is oscillation stable if this holds for
\(f(x)=\varphi(x,a)\) on every model and every infinite-dimensional subspace.

We use Hanson's Banach-language convention that the language contains
normalized linear-combination symbols \(\sigma_c\). Thus, for each scalar
tuple \(c=(c_1,\ldots,c_n)\), the pullback
\[
  (x_1,\ldots,x_n)\mapsto
  \varphi(\sigma_c(x_1,\ldots,x_n),a)
\]
is again a formula. Since \(T\) is stable, all such pullbacks are stable.

\section*{Referee Check and Repair}

The first version of this packet had two compressed steps. The
formula-aware Raynaud transfer was stated rather than unpacked, and the final
passage from type limits to an actual subspace was described as a subsequence
diagonalization. The latter wording is not strong enough: a property which
holds only for sufficiently far finite blocks does not automatically control
vectors involving the first few vectors of the closed span.

The proof below repairs this by using blocks, not a bare subsequence. Each new
block is chosen so that it realizes the same enriched type relative to the
finite net of the span already constructed. Thus every finite-dimensional
initial span is made almost monochromatic before passing to the closed span.

\section*{The Enriched Type Semigroup}

Fix a separable infinite-dimensional closed subspace \(X\subseteq Y\), a
countable rational vector subspace \(D\subseteq X\) which is dense in \(X\),
and a parameter \(a\). Let \(\Delta\) be a countable family of one-variable
formulas over \(D\cup\{a\}\) containing:
\begin{itemize}
\item all coordinates \(x\mapsto \|\lambda x+d\|\), with
  \(\lambda\in\Q\) and \(d\in D\);
\item the unit-ball coordinate \(x\mapsto \varphi(x,a)\), used only on
  normalized variables;
\item every formula obtained from the preceding coordinates by rational
  scalar dilation, rational finite sums, the normalized combination symbols
  \(\sigma_c\), and continuous connectives.
\end{itemize}
For formula coordinates whose output is evaluated on the unit ball, the
closure is taken through the normalized symbols \(\sigma_c\). The raw
convolution coordinates used to define the type semigroup are the norm-type
coordinates on sufficiently large Banach balls; the \(\varphi\)-coordinates
are read from the associated normalized pullbacks. We also close \(\Delta\)
under the two-variable pullbacks \(\theta(u+v)\) which occur when sums of
norm-type coordinates are formed. Since \(T\) is stable, every formula in this
fragment is stable.

For \(x\in D\), write
\[
  \bar x=(\theta(x))_{\theta\in\Delta}.
\]
Let \(\T\) be the pointwise closure of \(\{\bar x:x\in D\}\). We call its
elements \(\Delta\)-types. The norm coordinate
\(\nu(\tau)=\tau(\|x\|)\) defines bounded parts
\(\T_M=\{\tau:\nu(\tau)\le M\}\), and each \(\T_M\) is compact: the norm
coordinates put all representatives in a bounded ball of \(X\), while all
formula coordinates are uniformly bounded on bounded sorts.

\begin{lemma}[enriched convolution]\label{lem:enriched-conv}
For \(\tau,\rho\in\T\) and \(\alpha\in\Q\), choose defining sequences
\((x_i)\), \((z_j)\) in \(D\). The formulas
\[
  \alpha\tau=\lim_i\overline{\alpha x_i},
  \qquad
  \tau*\rho=\lim_i\lim_j\overline{x_i+z_j}
\]
define elements of \(\T\), independent of the defining sequences. Convolution
is commutative and associative, dilation distributes over convolution, and
for every rational tuple \(c=(c_1,\ldots,c_n)\) the map
\[
  \widehat T_c(\tau)=c_1\tau*\cdots*c_n\tau
\]
is Baire class one and is continuous at a comeager set of points in every
closed admissible conic class.
\end{lemma}

\begin{proof}
Fix a coordinate \(\theta\in\Delta\). For two defining sequences
\((x_i)\) and \((x'_i)\) of the same type and \((z_j)\), \((z'_j)\) of a
second type, the values of \(\theta(x_i+z_j)\) and
\(\theta(x'_i+z'_j)\) are governed by the stable formula \(\theta(u+v)\).
Grothendieck's double-limit criterion gives equality of the two possible
iterated limits, and also permits the order of the two limits to be
interchanged. This proves well-definedness and commutativity coordinate by
coordinate. Associativity and distributivity follow by applying the same
argument to the stable formula \(\theta(u+v+w)\) and to rational scalar
multiples.

For fixed \(\rho\), if \(\tau_m\to\tau\), choose realizations \(x_m\in D\)
with \(\bar x_m\) close to \(\tau_m\) and \(\bar x_m*\rho\) close to
\(\tau_m*\rho\). Then \((x_m)\) defines \(\tau\), so
\(\tau_m*\rho\to\tau*\rho\). Thus convolution is separately continuous.
On the compact annuli \(\T_{M+1}\setminus\operatorname{int}\T_M\), separately
continuous real coordinates are Baire class one; since the fragment is
countable, convolution and all countably many maps \(\widehat T_c\) have a
common comeager set of continuity points in each closed conic class.
\end{proof}

\section*{The Formula-Aware Raynaud Type}

A type is \emph{symmetric} if \((-1)\tau=\tau\). A nonzero closed set
\(\C\) of symmetric types is a \emph{conic class} if it is closed under
rational dilations and convolution. It is \emph{admissible} if it contains a
type with positive norm coordinate.

\begin{proposition}[formula-aware Raynaud type]\label{prop:raynaud}
There are \(1\le p<\infty\) and a symmetric admissible
\(\Delta\)-type \(\tau\), realized by blocks from \(X\), such that after
normalizing \(\nu(\tau)=1\),
\[
   \widehat T_c(\tau)=\|c\|_p\,\tau
   \tag{2}
\]
for every rational finite scalar tuple \(c\).
\end{proposition}

\begin{proof}
We give the exact Raynaud--Krivine--Maurey argument, noting at each point why
the added formula coordinates do not change it.

First, symmetric admissible types exist. Take a bounded separated sequence
\((u_i)\) in \(D\). Any cluster type \(q\) gives
\(q*(-q)\), the type of differences \(u_i-u_j\), and this type is symmetric
and has positive norm coordinate. Intersections of chains of closed
admissible conic classes remain admissible: scale an admissible element into
a fixed compact norm annulus and use the finite intersection property.
Therefore Zorn's lemma gives a minimal closed admissible conic class
\(\C\). By Lemma~\ref{lem:enriched-conv}, the countable family of maps
\(\widehat T_c\), \(c\in\Q^{<\N}\), has a common point of continuity
\(q\in\C\). Replacing \(q\) by a rational dilation if necessary, assume
\(\tau=q*(-q)\) is admissible.

Let \((x_i)\subseteq D\) define \(q\), and let
\(y_i=x_{2i-1}-x_{2i}\) define \(\tau\). Passing to a subsequence, form the
usual spreading model \(E\) of \((y_i)\), with spreading basis
\((e_i)\). A rational vector
\(\sum_{i=1}^n b_i e_i\) is said to realize
\(b_1\tau*\cdots*b_n\tau\). This realization convention is compatible with
convolution by construction.

The important comparison estimate is the following. If rational vectors
\(u,v\in\Span_\Q\{e_i\}\) realize enriched types \(\rho_u,\rho_v\), then for
every coordinate \(\theta\in\Delta\),
\[
  |\rho_u(\theta)-\rho_v(\theta)|
  \le \omega_\theta(\|u-v\|_E),
  \tag{3}
\]
where \(\omega_\theta\) is a modulus of uniform continuity for \(\theta\) on
a bounded sort containing the relevant representatives. Indeed, write \(u\)
and \(v\) with the same number of coordinates, realize them by far-out
blocks of the defining sequence \((y_i)\), apply the uniform continuity of
\(\theta\) in the original Banach norm, and pass to the iterated limits which
define both the spreading norm and the enriched types. This is the point at
which formula coordinates enter the proof. When \(\theta\) is a unit-ball
formula coordinate, it is first composed with the appropriate normalized
combination symbol \(\sigma_c\); hence all evaluations remain in the correct
sort, and the same uniform-continuity estimate applies.

For a rational tuple \(c=(c_1,\ldots,c_m)\), define the block operator
\[
  A_c e_k=\sum_{j=1}^m c_j e_{m(k-1)+j}.
\]
Then \(A_c u\) realizes \(\widehat T_c(\rho_u)\). The operators \(A_c\)
commute up to the tensor product rule
\[
  A_cA_d=A_{c\otimes d},
  \qquad
  \widehat T_c\widehat T_d=\widehat T_{c\otimes d}.
\]
The classical Krivine--Maurey approximate-eigenvector lemma for the
spreading basis says that for every finite commuting family
\(c^{(1)},\ldots,c^{(N)}\) there are numbers
\(\beta_1,\ldots,\beta_N\) and normalized rational vectors \(h_m\in E\) such
that
\[
  \|A_{c^{(\ell)}}h_m-\beta_\ell h_m\|_E\longrightarrow0
  \quad(1\le \ell\le N),
  \tag{4}
\]
with \(\|c^{(\ell)}\|_\infty\le\beta_\ell\le\|c^{(\ell)}\|_1\).
This is the finite joint approximate point-spectrum step in the original
Krivine--Maurey proof; it uses only the spreading basis and the Banach-space
norm, not the list of type coordinates.

Let \(\rho_m\) be the enriched type realized by \(h_m\). The set of all
limits of such approximating sequences is a closed admissible conic subclass
of \(\C\): dilation is immediate, convolution is obtained by putting two
realizing vectors on disjoint later supports, and closure follows by a
diagonal choice. Minimality of \(\C\) therefore implies that the common
continuity point \(q\), and hence \(\tau=q*(-q)\), is a limit of
approximating types satisfying (4).

Now evaluate at \(\tau\). Since \(q\) was chosen as a common continuity point
for all \(\widehat T_c\), and since (3) turns the norm convergence (4) into
coordinatewise convergence for every \(\theta\in\Delta\), we obtain
\[
  \widehat T_{c^{(\ell)}}(\tau)=\beta_\ell\,\tau
  \quad(1\le \ell\le N)
  \tag{5}
\]
as equality in the enriched type space. Thus every coordinate, including
\(\varphi(x,a)\), is carried along by the same equality.

Finally define \(N(c)\) by \(\widehat T_c(\tau)=N(c)\tau\). The norm
coordinate shows that \(N\) is a normalized, permutation-invariant,
1-unconditional norm on rational \(c_{00}\), satisfying
\[
  N(c\otimes d)=N(c)N(d),
  \qquad
  \|c\|_\infty\le N(c)\le\|c\|_1.
\]
The standard Krivine--Maurey multiplicative-norm calculation gives
\(N(c)=\|c\|_p\) for a single \(1\le p<\infty\). For completeness, recall
the calculation: \(N(1,\ldots,1)=n^{1/p}\) follows from tensor
multiplicativity and monotonicity by comparing \(n^k\) with powers of \(2\);
applying the same argument to tensor powers of an arbitrary positive rational
vector and using multinomial counting gives
\(N(c)^p=\sum_i |c_i|^p\); unconditionality gives signs. This proves (2).
\end{proof}

\section*{From The Type To A Real Subspace}

\begin{lemma}[block extraction]\label{lem:block-extraction}
Let \(\tau\) be as in Proposition~\ref{prop:raynaud}, normalized so that
\(\nu(\tau)=1\), and put \(r=\tau(\varphi(x,a))\). For every
\(\eps>0\) there is a normalized basic block sequence
\((z_k)\subseteq X\) such that for
\[
  Z=\overline{\Span}\{z_k:k\in\N\}
\]
one has
\[
  |\varphi(u,a)-r|<\eps
  \qquad (u\in S_Z).
\]
\end{lemma}

\begin{proof}
Let \((y_i)\subseteq X\) be a defining block sequence for \(\tau\). Equality
(2) says that if \(c\in\Q^n\) and \(\|c\|_p=1\), then the enriched type of
the corresponding normalized block combination, in the iterated block limit,
is again \(\tau\). In particular, both its norm coordinate and the
\(\varphi(\sigma_c(\bar x),a)\)-coordinate are the same as those of \(\tau\).

We construct \(z_1,z_2,\ldots\) recursively as finite rational blocks of
\((y_i)\), with successive disjoint supports. Fix an increasing exhaustion
\(\Delta_1\subseteq\Delta_2\subseteq\cdots\) of the fragment by finite sets
of coordinates and choose \(\delta_m\downarrow0\). The induction invariant is
stronger than small oscillation: for every rational unit vector
\(b=(b_1,\ldots,b_m)\) in a prescribed finite \(\delta_m\)-net of the
\(\ell_p^m\)-sphere, the normalized block vector
\(\sigma_b(z_1,\ldots,z_m)\) has \(\Delta_m\)-type within \(\delta_m\) of
\(\tau\).

Assume \(z_1,\ldots,z_{m-1}\) have been chosen. Take a finite rational net
for the unit sphere of \(\ell_p^m\). For a net vector
\((b_1,\ldots,b_m)\), put
\(s=\left(\sum_{j<m}|b_j|^p\right)^{1/p}\). By the induction hypothesis,
when \(s>0\) the normalized previous part
\(s^{-1}\sum_{j<m}b_jz_j\) has enriched type close to \(\tau\) on a finite
coordinate set. Equality (2) gives
\[
  s\tau*b_m\tau=(s^p+|b_m|^p)^{1/p}\tau=\tau
\]
for the corresponding model combination. Hence the requirements that
\(\sigma_b(z_1,\ldots,z_m)\) be \(\Delta_m\)-close to \(\tau\), for all net
vectors \(b\), form a finite set of open conditions in the enriched type
topology. Since \(\tau\) is realized by arbitrarily far normalized blocks of
\((y_i)\), a far-out rational block can be chosen as \(z_m\) satisfying all
these conditions. The case \(s=0\) is just the requirement that \(z_m\)
itself realize \(\tau\) on \(\Delta_m\).

This recursive construction has the following consequence. For each \(m\),
every unit vector in the finite-dimensional space
\(\Span(z_1,\ldots,z_m)\) has \(\varphi\)-value within \(\eps/2\) of \(r\):
approximate it by a point of the rational net used at stage \(m\), then use
the common modulus of uniform continuity of \(\varphi\) and the fact that the
net point has \(\Delta_m\)-type close to \(\tau\). The norm coordinates in
the same finite sets make \((z_k)\) a basic sequence, because the
finite-dimensional norms are uniformly close to the corresponding
\(\ell_p\)-sum norms on the chosen nets.

Now let \(u\in S_Z\). Pick \(m\) and
\(u_m\in S_{\Span(z_1,\ldots,z_m)}\) with \(\|u-u_m\|\) so small that
uniform continuity changes \(\varphi\) by less than \(\eps/2\). The previous
paragraph gives \(|\varphi(u_m,a)-r|<\eps/2\), hence
\(|\varphi(u,a)-r|<\eps\).
\end{proof}

\section*{Main Theorem}

\begin{theorem}\label{thm:main}
Let \(T\) be a stable Banach theory. Then every unary formula of \(T\) is
oscillation stable on models of \(T\).
\end{theorem}

\begin{proof}
Let \(M\models T\), let \(\varphi(x,a)\) be a unary formula instance, let
\(Y\subseteq M\) be an infinite-dimensional closed subspace, and fix
\(\eps>0\).

Choose a separable infinite-dimensional closed subspace \(X\subseteq Y\), and
build the enriched fragment over \(X\) and \(a\) as above. By
Proposition~\ref{prop:raynaud} and Lemma~\ref{lem:block-extraction}, there is
an infinite-dimensional closed subspace \(Z\subseteq X\subseteq Y\) and a
real number \(r\) such that
\[
  |\varphi(u,a)-r|<\eps/2
  \qquad (u\in S_Z).
\]
Therefore, for all \(u,v\in S_Z\),
\[
  |\varphi(u,a)-\varphi(v,a)|<\eps
\]
so \(\varphi(x,a)\) is oscillation stable on \(Y\). Since \(M,Y,a\), and
\(\eps\) were arbitrary, the theorem follows.
\end{proof}

\section*{Why The Naive Symmetry Route Is Not Used}

The proof does not rely on the false principle that a symmetric uniformly
continuous coloring on the Hilbert sphere becomes constant on vectors with
small \(\ell_\infty\)-coordinates. That principle fails: if \(u_n\) is the
flat vector of support \(n\) in \(\ell_2\), and \(A\) is the union of the
sign/permutation orbits of \(u_{2^k}\), then \(x\mapsto d(x,A)\) is
Lipschitz and symmetric but separates the support sizes \(2^k\) and
\(3\cdot2^k\) by a fixed amount.

The present argument uses more than symmetry. It uses stability of every
normalized linear-combination pullback, which is precisely what makes the
formula coordinate part of the Krivine--Maurey type equality.

\section*{Verification Notes}

The proof was checked against three prior failed routes:
\begin{itemize}
\item The compactness/Dvoretzky--Milman Hilbert-subspace route fails because
  finite good subspaces need not be nested.
\item The minimal-wide-type route isolates the same obstruction as preservation
  of tail finite satisfiability through the Hilbert Morley spine.
\item The bare symmetric-flattening lemma is false, as explained above.
\end{itemize}

The formula-aware Krivine--Maurey route avoids all three. It builds the
sequence inside the original subspace \(Y\), and the formula value is a
coordinate of the same stable type equality that produces the \(\ell_p\)
spreading behavior.

Human review should focus on Proposition~\ref{prop:raynaud} and
Lemma~\ref{lem:block-extraction}. Proposition~\ref{prop:raynaud} is where the
exact Raynaud/Krivine--Maurey conic-class proof is transferred from norm types
to the enriched stable type space. Lemma~\ref{lem:block-extraction} is the
repair of the old subsequence diagonalization: it uses blocks so that finite
initial spans, and hence the closed span, are controlled.

\section*{Novelty Check}

Local run indexes were searched for arXiv:2004.03062, Hanson, Question 4.3,
minimal wide types, oscillation stability, Krivine--Maurey, and stable Banach
theory. Existing durable work contained only a Hilbert-operator partial result,
a WAP conditional reduction, and failed attempt notes.

Fresh web searches on 2026-06-16 for exact and close phrases including
\begin{itemize}
\item ``Does (model theoretic) stability imply oscillation stability''
\item ``Hanson oscillation stability stable Banach theory''
\item ``stability implies oscillation stability Banach unary formula''
\item ``Question 4.3 Indiscernible Subspaces oscillation stability''
\end{itemize}
found Hanson's source paper, the absorbing paper, and the
Shelah--Usvyatsov minimal-types paper. They did not reveal a later explicit
answer to Question 4.3. The proof here appears to be a new synthesis of
Hanson's question with the Krivine--Maurey/Raynaud stable type machinery.

\section*{References}

\begin{thebibliography}{9}

\bibitem{Hanson2020}
James Hanson.
\newblock \emph{Indiscernible Subspaces and Minimal Wide Types}.
\newblock arXiv:2004.03062.

\bibitem{Hanson2020b}
James Hanson.
\newblock \emph{Strongly Minimal Sets and Categoricity in Continuous Logic}.
\newblock arXiv:2011.00610.

\bibitem{KM1981}
Jean-Louis Krivine and Bernard Maurey.
\newblock \emph{Espaces de Banach stables}.
\newblock Israel Journal of Mathematics 39 (1981), 273--295.

\bibitem{Raynaud1983}
Yves Raynaud.
\newblock \emph{Espaces de Banach superstables, distances stables et
homeomorphismes uniformes}.
\newblock Israel Journal of Mathematics 44 (1983), 33--52.

\bibitem{BenYaacov2013}
Itai Ben Yaacov.
\newblock \emph{Model theoretic stability and definability of types, after
A. Grothendieck}.
\newblock arXiv:1306.5852.

\bibitem{ShelahUsvyatsov2014}
Saharon Shelah and Alexander Usvyatsov.
\newblock \emph{Minimal types in stable Banach spaces}.
\newblock arXiv:1402.6513.

\bibitem{BragaSwift2017}
Bruno M. Braga and Christopher Swift.
\newblock \emph{Coarse and uniform embeddings into superstable Banach spaces}.
\newblock arXiv:1704.04468.

\end{thebibliography}

\end{document}
